\chapter{Implementacja HTML i CSS}

\section{Zakres wdrożenia i responsywność}
Zgodnie z wymogami projektu, wybrane widoki zaprojektowane w programie Figma zostały zaimplementowane do działającego kodu przy użyciu technologii HTML5 oraz CSS3. 

Kluczowym założeniem podczas kodowania było zapewnienie pełnej responsywności (RWD - Responsive Web Design). Dzięki zastosowaniu nowoczesnych technik ostylowania (takich jak Flexbox oraz CSS Grid) a także zapytań medialnych (Media Queries), zaimplementowane interfejsy płynnie skalują się i dostosowują swój układ do szerokości okna przeglądarki. Potwierdza to, że system jest w pełni gotowy do obsługi zarówno na urządzeniach desktopowych, jak i mobilnych (smartfony, tablety).

Zgodnie z wymogiem prowadzącego, w procesie tworzenia bazowej struktury HTML oraz podstawowych reguł CSS wykorzystano wsparcie narzędzi AI (np. ChatGPT), przy czym wygenerowany kod został ręcznie zoptymalizowany, przetestowany i dostosowany do rygorystycznych wymogów responsywności.

\section{Struktura katalogów}
Aby zachować maksymalną czytelność i porządek w kodzie źródłowym, projekt został podzielony na logiczne foldery odpowiadające wdrożonym widokom. Poniżej przedstawiono strukturę katalogów zaimplementowanej części interfejsu:

\begin{itemize}
    \item \textbf{MainPage} -- Strona główna z ofertą.
    \begin{itemize}
        \item \texttt{MainPage.html} -- semantyczna struktura widoku
    \end{itemize}
    
    \item \textbf{Dashboard} -- Panel użytkownika z historią zamówień.
    \begin{itemize}
        \item \texttt{dashboard.html} -- struktura widoku panelu klienta
    \end{itemize}

    \item \textbf{OrderConfigurator} -- Zaawansowany konfigurator zamówienia.
    \begin{itemize}
        \item \texttt{configurator.html} -- struktura formularzy i kalendarza
        \item \texttt{style.css} -- style formularzy i siatki (Grid/Flexbox)
        \item \texttt{script.js} -- plik obsługujący interakcję (np. wybór dat)
    \end{itemize}

    \item \textbf{AdminDiets-menu} -- Moduł zarządzania dietami dla administratora.
    \begin{itemize}
        \item \texttt{Diets-menu-managment.html} -- struktura widoku panelu administracyjnego
    \end{itemize}
\end{itemize}

Takie podejście ułatwia ewentualną dalszą rozbudowę aplikacji oraz integrację front-endu z systemami back-endowymi.